\documentclass[12pt]{article}
\usepackage[utf8]{inputenc}
\usepackage[T1]{fontenc}
\usepackage[margin=1in]{geometry}
\usepackage{setspace}
\usepackage{xcolor}
\usepackage[colorlinks=true, urlcolor=blue!60!black, linkcolor=blue!60!black]{hyperref}
\setlength{\parindent}{0em}
\setlength{\parskip}{0.45em}
\renewcommand{\baselinestretch}{1.0}
\begin{document}
\begin{center}
{\Large \textit{Dolan v. City of Tigard}, 512 U.S. 374 (1994)\par}
\vspace{0.3em}
{\normalsize Case Brief --- One-Page Edition\par}
\end{center}

\textbf{Memory Jogger.} Permit conditions requiring land dedications must be roughly proportional to the project's impact, as shown by an individualized determination (\href{https://supreme.justia.com/cases/federal/us/512/374/}{Justia}).

\textbf{Facts.} Florence Dolan sought to nearly double her plumbing and electric supply store beside Fanno Creek in Tigard, Oregon. The city required her to dedicate the floodplain portion as a public greenway and a 15-foot strip for a public bike path---about 10\% of her parcel---finding that the path ``could offset some of the traffic demand.''

\textbf{Procedural History.} Oregon's land-use board and courts upheld the conditions as reasonably related to the expansion.

\textbf{Votes.} Reversed, 5--4: Rehnquist for the Court; Stevens (with Blackmun and Ginsburg) and Souter dissented.

\textbf{Holding.} Beyond \textit{Nollan}'s essential nexus, the city must show that a dedication is related in nature and extent to the development's impact; neither condition did.

\textbf{Reasoning.} Demanding the land outright would be a per se taking, so the unconstitutional conditions doctrine governs. The conditions related to flood control and traffic, but the city never explained why flood control needed public access, and a path that ``could'' offset traffic was not shown to do so. The burden lies on the city. Stevens saw a revival of \textit{Lochner}; Souter called the defects nexus failures. \emph{Later:} \textit{Koontz} (2013) reached money demands and \textit{Sheetz} (2024) legislative fees.

\textbf{Examples.} \emph{Same-Side:} a \$60,000 traffic fee on a backyard cottage. \emph{Opposite-Side:} a turn lane backed by a traffic study. \emph{Fence-Sitter:} a citywide school fee on a retirement community.

\textbf{Critique.} The burden shift may push cities to deny permits outright or rely on broad taxes instead of tailored fees.

\textbf{Key Quotation.} ``We think a term such as `rough proportionality' best encapsulates what we hold to be the requirement of the Fifth Amendment. No precise mathematical calculation is required, but the city must make some sort of individualized determination that the required dedication is related both in nature and extent to the impact of the proposed development.'' --- Chief Justice Rehnquist, 512 U.S. at 391.

\end{document}
